\documentclass[11pt,a4paper]{article}

\usepackage[margin=1in]{geometry}
\usepackage{amsmath,amssymb,amsthm}
\usepackage{booktabs}
\usepackage{array}
\usepackage{listings}
\usepackage{xcolor}
\usepackage{framed}
\usepackage{eurosym}
\usepackage[authoryear,round]{natbib}
\usepackage[hidelinks]{hyperref}
\usepackage{microtype}

\lstdefinelanguage{yaml}{
  keywords={true,false,null},
  sensitive=false,
  comment=[l]{\#},
  morestring=[b]',
  morestring=[b]"
}

\lstset{
  basicstyle=\ttfamily\small,
  commentstyle=\color{gray},
  stringstyle=\color{black},
  breaklines=true,
  frame=single,
  framesep=4pt,
  numbers=left,
  numberstyle=\tiny\color{gray},
  captionpos=b
}

\newtheorem{theorem}{Theorem}[section]
\newtheorem{corollary}{Corollary}[section]
\newtheorem{principle}{Principle}[section]
\theoremstyle{definition}
\newtheorem{definition}{Definition}[section]
\theoremstyle{remark}
\newtheorem{remark}{Remark}[section]

\title{Boundary Discipline for Accounting Information Systems:\\
A Design-Science Architecture for Preserving Verification Failure}

\author{}
\date{}

\begin{document}

% --------------------------------------------------------------------
% Title page (separate page, author-identified). Detach this page before
% sending the manuscript out for blind review, per IJDAR's anonymous
% review policy; the manuscript body below carries no author identification.
% --------------------------------------------------------------------
\begin{titlepage}
\centering
\vspace*{2cm}
{\LARGE\bfseries Boundary Discipline for Accounting Information Systems:\\
A Design-Science Architecture for Preserving Verification Failure\par}
\vspace{2cm}
{\large Duston Moore, PhD\par}
\vspace{0.5em}
{\large Independent Researcher\par}
{\large Edmonton, Alberta, Canada\par}
\vspace{0.5em}
{\large Corresponding author e-mail: \texttt{dhwcmoore@gmail.com}\par}
\vspace{2cm}
% NOTE: set this to the actual month of submission before sending;
% left unset rather than carrying forward a date that has already
% passed and would need to be caught and corrected again.
{\large [MONTH YEAR -- set at submission]\par}
\thispagestyle{empty}
\end{titlepage}

\newpage
\pagenumbering{arabic}

\begin{abstract}
Auditing and accounting information systems increasingly rely on automated classification, and these systems routinely encounter cases in which evidence is missing, jurisdictionally obstructed, or inconclusive. The dangerous failure mode is not uncertainty itself but silent conversion: the operational treatment of an unresolved assertion as accepted, without an explicit indeterminate, modified, blocked, or escalated classification, because the workflow demands a binary output. This paper asks how an accounting information system should be designed so that verification failure remains a durable, consequence-bearing state from evidence intake through audit reporting, residual-risk assignment, inspection, and digital disclosure, and answers it with a design-science architecture, Boundary Discipline, constructed and evaluated against explicit criteria rather than argued for a priori. Its components are explicit verification boundary specifications; a multi-valued classification scheme in which \textsf{Undefined} is a success state, not an error condition; a Verification Balance Sheet developed here as an accounting instrument, with required fields, responsible organisational roles, and creation and review points, not only a formal definition; a Risk Assignment Register developed as a durable organisational record connecting each residual to a proposed owner, that owner's accepted capacity, orphan status, escalation authority, and disclosure status; an organisational responsibility model separating evidence provision, proposal, classification, admission, ownership, capacity review, and escalation into distinct roles, with registry-level identity separation stated precisely as what it does and does not establish about organisational independence; a composition rule under which any single unresolved material dependency blocks an opinion regardless of how many others are verified; a standards-mapping analysis against ISA 500, 505, 580, and 705 and their national analogues; precisely defined inspection measures for boundary coverage, undefined-region materiality, residual assignment, and orphan rate; and a proposed machine-readable digital-disclosure extension for XBRL and ESEF filings. The Wirecard collapse, corrected against the documented history -- the Philippine banks were reachable and responded promptly once asked directly; the failure was an uncontrolled trustee-mediated confirmation route, not jurisdiction -- serves as the running case, connected explicitly to the balance sheet, the register, and audit-committee escalation, and the silent absorption of deprioritised exceptions in continuous auditing serves as the ordinary-scale analogue. Selected classification, admission, and residual-preservation requirements this architecture depends on are additionally validated, in one sentence: a companion Coq/Rocq kernel and its extracted OCaml classifier demonstrate that these requirements can be implemented without an affirmative default, which is supporting evaluation evidence for this architecture, not this paper's own technical contribution.
\end{abstract}

\noindent\textbf{Keywords:} accounting information systems; continuous auditing; audit evidence; classification integrity; formal verification

\medskip
% NOTE: this repository link must point to an anonymised mirror in the
% review copy -- this journal's anonymous review policy is confirmed
% (see the titlepage NOTE above): generate the anonymised link at
% https://anonymous.4open.science/ (paste the repository URL, add the
% author's name/username/email as terms to hide) and insert it below
% before submission. Replace with the de-anonymised repository URL
% and the specific commit or tag only in the camera-ready version.
\noindent\textbf{Data availability statement:} This is a conceptual and methodological paper. It does not analyse a proprietary or empirical dataset. The illustrative boundary specification and classification pseudocode presented in the paper are original artefacts of the manuscript and are reproduced here in full. The classification, admission, and residual-preservation requirements Section~\ref{sec:formal} identifies as selected for formal validation (Table~\ref{tab:validation}) are additionally formalised and machine-checked, generally rather than only for the illustrative cases of Sections~\ref{sec:wirecard} and~\ref{sec:ca-case}, as supporting evaluation evidence rather than as this paper's own technical contribution, in a companion Coq/Rocq development and extracted OCaml classifier -- itself the companion manuscript's central contribution -- at an anonymised mirror of the repository, \url{[ANONYMISED ARTEFACT LINK -- insert before submission]}; the de-anonymised repository and the specific commit or tag corresponding to this version of the manuscript will be recorded here in the camera-ready version, once the artefact is finalised for submission. Section~\ref{sec:limitations} states what this development does and does not establish; the companion repository's \texttt{NON\_CLAIMS.md} gives the complete, itemised list.

% NOTE: "will be disclosed... before either submission" below is
% deliberately future tense -- do not change to past tense until the
% actual notification has been sent and its date can be recorded here.
\medskip
\noindent\textbf{Note on a related submission:} A companion manuscript, \emph{A Discipline of Partial Audit Classification: A Machine-Checked Kernel for Default-Free Audit Decisions}, is being prepared for submission to the \emph{International Journal on Software Tools for Technology Transfer}. The manuscripts share their motivating problem, central conceptual framework, classification vocabulary, and Coq/Rocq artefact -- part of the conceptual contribution is shared, and this note says so rather than describing the overlap as only a shared invariant and proof artefact. The present manuscript's exclusive primary contribution is the accounting-information-systems architecture: its design-science formulation, Verification Balance Sheet, Risk Assignment Register, organisational responsibility model, standards mapping, inspection measures, and digital-reporting implications. The companion manuscript's exclusive primary contribution is the executable kernel: its formal semantics, decision procedures, certificate validator, residual transition system, theorem suite, extraction, trusted-computing-base analysis, and implementation evaluation. Both complete manuscripts will be disclosed to both editors before either submission, so that the extent and acceptability of the overlap can be assessed directly. Neither manuscript claims that the shared material is exclusive to it.

\section{Introduction: Verification Failure and Silent Conversion}
\label{sec:intro}

Audit systems and accounting information systems classify claims. A management assertion about the existence of a cash balance, a transaction flagged by a continuous monitoring rule, and a field populated in a risk dashboard are all inputs to classification processes whose outputs bind real consequences: unqualified opinions, cleared exception queues, and assurances passed upward to boards and markets.

Classification outside verified domains is dangerous, but the danger is not totality itself. A classifier forced to produce an affirmative or otherwise determinate verdict for every input, including inputs over which it has no verified authority, will assign arbitrary or convenient categories precisely where restraint is required; in audit settings, the convenient category is acceptance. A default-free classifier may remain total by returning \textsf{Undefined} outside its verified domain -- this paper's own classifier does exactly that, for every input, without ever becoming the dangerous kind. The error is therefore not totality itself, but total \emph{affirmative} or determinate-and-binding classification without total verified coverage. When verification procedures fail, whether through jurisdictional obstruction, missing third-party data, ambiguous control relationships, or inconclusive automated checks, the surrounding workflow frequently proceeds as though the assertion had been verified. We call this failure mode \emph{silent conversion}: the operational transformation of verification failure into audit acceptance, without an explicit indeterminate, modified, blocked, or escalated classification, and often without any durable trace in the audit trail. The failure has a precise form. It is a type error: an assertion of type \textsf{Undefined} is used where an assertion of type \textsf{Verified} is required, and the remainder of this paper is the discipline that makes the type error inadmissible.

The collapse of Wirecard AG in June 2020 illustrates silent conversion at catastrophic scale \citep{mccrum2022, esma2020}. Ernst \& Young audited Wirecard's financial statements for roughly a decade, from 2009 until the 2020 collapse, issuing unqualified opinions throughout; the specific \officialeuro 1.9 billion figure is what stood, in the 2019 financial statements under review in June 2020, as the balance purportedly held in trustee accounts at two Philippine banks, not a constant amount asserted across the full decade of the engagement. For years, confirmation of those balances was obtained not directly from the banks but through a Manila-based third-party trustee, who supplied letters and documents on what appeared to be the banks' letterhead; no one at the auditor is recorded as having verified the accounts by direct contact with the banks themselves during that period. When direct confirmation was finally sought in June 2020, both banks responded -- promptly, and through the jurisdiction's normal channels -- that no such accounts had ever existed and that the supporting documents were fabricated. The conventional explanations, inadequate professional scepticism or sophisticated fraud, are insufficient. The deeper problem is structural: the operative audit methodology permitted a confirmation route the auditor did not control, and evidence of a grade lower than direct bank confirmation, to be treated as equivalent to verification success, for years, before the jurisdiction was ever shown to be an obstacle. When direct bank confirmation was not obtained, alternative documentation routed through an intermediary was accepted as sufficient in its place. The methodology converted cannot verify into verified. This account is independently corroborated by the special investigation conducted for the German Bundestag's Wirecard inquiry committee \citep{wambach2021}, which examined Ernst \& Young's own working papers and reported the same substitution of intermediary-routed assurance for primary confirmation on the relevant balances; \citet{beerbaum2020} discusses that investigation's findings but is not itself the investigation.

Silent conversion is not confined to spectacular frauds. It occurs in ordinary accounting information systems whenever a continuous auditing dashboard clears deprioritised exceptions that were never investigated, whenever a missing third-party confirmation becomes a no exception noted entry, and whenever a null risk field is treated downstream as zero risk. In each case an epistemic state, namely unresolved, is collapsed into an operational state, namely accepted, because the workflow needs a binary output.

This paper asks a single research question, and it governs everything that follows: how should an accounting information system be designed so that verification failure remains a durable, consequence-bearing state from evidence intake through audit reporting, residual-risk assignment, inspection, and digital disclosure?

The answer developed here is a methodology we term \emph{Boundary Discipline}. This paper's contribution is nine pieces built as one design-science artefact, not nine separate claims:

\begin{enumerate}
\item silent conversion identified as an accounting information systems failure pattern that recurs across evidence intake, application logic, exception workflows, and audit consequence, not a failure confined to any one of them (Section~\ref{sec:intro}, Section~\ref{sec:lit});
\item Boundary Discipline itself, a design-science accounting information systems architecture built and evaluated against explicit criteria (Section~\ref{sec:method}, Section~\ref{sec:architecture});
\item the Verification Balance Sheet, an accounting instrument delimiting what a verification procedure's target, assumptions, cost, residuals, and trust chain actually are (Section~\ref{sec:vbs});
\item the Risk Assignment Register, a durable organisational record connecting every residual to an owner, that owner's capacity, orphan status, escalation authority, and disclosure (Section~\ref{sec:rar});
\item an organisational responsibility model separating proposal, classification, admission, ownership, review, and escalation into distinct roles (Section~\ref{sec:roles});
\item a standards-mapping analysis across ISA 500, ISA 505, ISA 580, and ISA 705, their national analogues, and digital-reporting requirements, identifying what each already requires and what Boundary Discipline operationalises beyond it (Section~\ref{sec:implications});
\item inspectable, precisely defined measures for boundary coverage, undefined-region materiality, residual assignment, and orphan rate (Section~\ref{sec:inspection});
\item demonstration through the historically corrected Wirecard case and a continuous-auditing workflow, each connected explicitly to the Verification Balance Sheet, the Risk Assignment Register, and the responsibility model (Section~\ref{sec:wirecard}, Section~\ref{sec:ca-case});
\item formal validation, through a companion Coq/Rocq kernel and its executable extraction, that selected classification, admission, and preservation requirements this architecture depends on can be implemented without an affirmative default -- supporting evidence for the architecture, not this paper's own technical contribution (Section~\ref{sec:coq}).
\end{enumerate}

One rule governs everything that follows, and it is stated before any component is described.

\begin{principle}[Central Invariant]
\label{prin:invariant}
No unmodified audit opinion may rely affirmatively on a material assertion unless that assertion is classified Verified and validly admitted. Any material assertion that is unresolved must constrain, modify, or preclude the opinion -- it may not simply be absent from the reasoning that produced an unqualified conclusion. A qualified opinion or disclaimer may, and typically does, depend on exactly such an unresolved assertion; the invariant restricts what an \emph{unqualified} opinion may depend on, not what any opinion may mention.
\end{principle}

In symbolic form:

\begin{lstlisting}[numbers=none]
Unqualified(Opinion) =>
  for every material assertion a used by Opinion:
      state(a) = Verified and admitted(a)
\end{lstlisting}

Equivalently, and more sharply:

\begin{lstlisting}[numbers=none]
state(a) = Undefined => a is inadmissible for an unqualified opinion
\end{lstlisting}

Every component of the methodology is support for this invariant. The boundary specifications compute the classification states, the residual register preserves the non-affirmative ones, the composition rule propagates them upward, and the admissibility theorem of Section~\ref{sec:formal} excludes any unqualified conclusion that depends on an \textsf{Undefined} material assertion.

The paper proceeds as follows. Section~\ref{sec:lit} reviews how audit theory, accounting information systems research, database systems, and runtime verification each address part of the verification failure problem, and identifies the gap that Boundary Discipline fills. Section~\ref{sec:concepts} defines the core concepts. Section~\ref{sec:principles} states design principles for default-free audit classification. Section~\ref{sec:architecture} presents Boundary Discipline as an accounting information systems architecture. Section~\ref{sec:formal} gives the formal specification and proof-of-concept. Sections~\ref{sec:wirecard} and~\ref{sec:ca-case} apply the framework to the Wirecard case and to continuous auditing exceptions respectively. Section~\ref{sec:implications} discusses implications for standards and inspection. Section~\ref{sec:limitations} states limitations and future work, and Section~\ref{sec:conclusion} concludes.

\subsection{Research Method}
\label{sec:method}

This paper follows the design-science research paradigm for information systems: it treats Boundary Discipline as an artefact constructed to address an identified problem, evaluated against explicit criteria, rather than as a position advanced by argument alone \citep{hevner2004, peffers2007}. The six activities of that paradigm map onto this paper's sections as follows, and Table~\ref{tab:dsrmap} restates the mapping at the level of individual failure modes, mechanisms, and evidence.

\begin{enumerate}
\item \textbf{Problem identification} (Section~\ref{sec:intro}, Section~\ref{sec:lit}): silent conversion, motivated by the Wirecard collapse and grounded against what audit theory, AIS research, database semantics, and runtime verification each already supply and each leave unaddressed.
\item \textbf{Design objectives} (Section~\ref{sec:principles}): the six principles a default-free classification system must satisfy, derived from the failure modes identified in Section~\ref{sec:lit}.
\item \textbf{Artefact construction} (Section~\ref{sec:concepts}, Section~\ref{sec:architecture}): the Verification Balance Sheet, the Risk Assignment Register, and the classification pipeline -- boundary specification, admission certificate, residual register, composition rule.
\item \textbf{Formal verification} (Section~\ref{sec:formal}): the design objectives restated as machine-checked theorems over an executable implementation of the artefact, not only as specification.
\item \textbf{Demonstration} (Sections~\ref{sec:wirecard}--\ref{sec:ca-case}): the artefact applied to the Wirecard collapse and to a continuous-auditing exception queue.
\item \textbf{Evaluation and limitations} (Section~\ref{sec:limitations}): six evaluation activities, none of them empirical -- requirements traceability against Section~\ref{sec:lit}'s identified gaps, internal logical consistency among the principles of Section~\ref{sec:principles}, formal validation of the selected implementation properties Table~\ref{tab:validation} reports, historical case reconstruction (Section~\ref{sec:wirecard}), continuous-auditing workflow demonstration (Section~\ref{sec:ca-case}), and boundary and non-claim analysis stated as narrowly as Section~\ref{sec:coq}'s formal claims are. Section~\ref{sec:limitations} also states the future evaluation programme this paper does not itself carry out.
\end{enumerate}

\begin{table}[htbp]
\centering
\small
\begin{tabular}{>{\raggedright\arraybackslash}p{2.8cm}>{\raggedright\arraybackslash}p{3.2cm}>{\raggedright\arraybackslash}p{3.2cm}>{\raggedright\arraybackslash}p{3.4cm}}
\toprule
\textbf{Failure mode} & \textbf{Architectural mechanism} & \textbf{Formal validation (Table~\ref{tab:validation})} & \textbf{Demonstration} \\
\midrule
Unverified treated as accepted & Boundary specification, \textsf{Undefined} as a first-class output & No affirmative default & Wirecard: confirmation-route failure (Section~\ref{sec:wirecard}) \\
\addlinespace
Proposer admits its own claim & Certificate-based admission against a process registry & Admission separation & Six adversarial fixtures in the companion development (Section~\ref{sec:coq}) \\
\addlinespace
Deprioritised exceptions vanish from record & Append-only Risk Assignment Register & Residual preservation & Continuous auditing: dashboard clearing (Section~\ref{sec:ca-case}) \\
\addlinespace
Partial evidence quietly carries an opinion & Composition: any unresolved material dependency blocks the opinion & Opinion restriction & Both cases: the opinion remains inadmissible under partial evidence \\
\bottomrule
\end{tabular}
\caption{Failure modes, the mechanism each is answered by, the formally validated kernel property that supports it, and where it is demonstrated.}
\label{tab:dsrmap}
\end{table}

\subsection*{Core Terms}

\begin{framed}
\noindent\textbf{Verification failure} occurs when a system cannot determine whether a claim is verified or refuted within the evidence, authority, and observation boundaries available to it.

\medskip
\noindent\textbf{Silent conversion} occurs when a verification failure is operationally treated as acceptance without an explicit \textsf{Undefined}, modified, blocked, or escalated classification.

\medskip
\noindent\textbf{Implicit default} is the local rule, convention, or implementation path by which silent conversion occurs.

\medskip
\noindent\textbf{Boundary} is the demarcation between regions where a verification procedure produces valid results and regions where it does not.

\medskip
\noindent\textbf{Undefined} is the classification result when an input falls outside all defined verification boundaries.

\medskip
\noindent\textbf{Residual} is what remains unverified after a verification procedure completes, including both known limitations and structural blind spots.

\medskip
\noindent\textbf{Risk assignment} is the process of identifying who bears responsibility for losses if a verification residual materialises as actual error.

\medskip
\noindent\textbf{Admission} is the act by which a distinct, registered verifier identity accepts a proposed classification as binding for audit consequence.
\end{framed}

\section{Verification Failure Across Auditing, AIS, and Formal Methods}
\label{sec:lit}

The problem addressed in this paper is not newly invented. Each of the literatures reviewed below has confronted verification failure within its own domain and has developed partial responses. What is missing is a methodology that joins these responses into a single classification architecture for audit information systems.

\subsection{Audit Evidence, Professional Scepticism, and Verification Failure}
\label{sec:lit-audit}

Audit theory has long recognised that evidential insufficiency is a distinct epistemic condition rather than a weak form of assurance. Mautz and Sharaf's foundational treatment of auditing as an evidence-gathering discipline distinguishes the propositions an auditor can establish from the propositions an auditor merely fails to refute \citep{mautz1961}. The audit quality and audit risk literature likewise treats the sufficiency and appropriateness of evidence as the central determinant of the assurance an engagement can support \citep{knechel2013}. The professional scepticism literature develops this distinction behaviourally, characterising scepticism as a disposition to withhold conclusion until sufficient appropriate evidence is obtained rather than as generalised distrust \citep{nelson2009, hurtt2010}.

The standards themselves encode a partial recognition of verification failure. Under ISA 705 and its national analogues, an inability to obtain sufficient appropriate audit evidence constitutes a scope limitation, and a scope limitation affecting material items requires a qualified opinion or a disclaimer of opinion \citep{isa705}. The national analogues are not merely formal restatements: Canada's CAS 705 and the equivalent PCAOB standard governing departures from an unqualified opinion, AS 3105, encode the same qualified-versus-disclaimer distinction as a function of whether the possible effects of undetected misstatement are material but not pervasive, or both material and pervasive \citep{cas705, pcaob3105}, which is precisely the materiality-conditioned branching that the boundary specification in Section~\ref{sec:formal} renders machine-checkable. The empirical literature on modified opinion decisions shows that the modification channel operates systematically in practice, though its use is shaped by client characteristics and auditor incentives rather than governed by any classification architecture \citep{habib2013}. The logical structure of the modified opinion regime is instructive: it acknowledges that failed verification is not verification, and it provides an output channel, the modified opinion, that is distinct from acceptance.

Two weaknesses limit this recognition in practice. First, the standards permit alternative evidence pathways whose epistemic grade is markedly lower than the primary procedure they replace. ISA 500 requires audit evidence to be sufficient and appropriate without ranking evidence sources against one another \citep{isa500}; ISA 505 treats external confirmation, sought and controlled by the auditor, as ordinarily more reliable evidence for existence assertions specifically \citep{isa505}; and ISA 580 is explicit that written representations from management, on their own, do not provide sufficient appropriate audit evidence about the matters they address \citep{isa580} -- a standard directly on point for the Wirecard pattern of Section~\ref{sec:wirecard}, where representations and intermediary-routed documentation stood in for a confirmation the auditor never controlled. The evidence-grade hierarchy these three standards jointly imply is exactly the ranking, primary confirmation above tertiary representation, that Table~\ref{tab:boundaryspec}'s boundary specification makes mechanical; what the standards recognise as a matter of professional judgement, this architecture recognises as a matter of what a procedure's postcondition is entitled to certify. The distinction between absence of evidence and evidence of absence is preserved in principle but eroded in the operative decision procedure. Second, the recognition is behavioural and documentary rather than architectural. Nothing in the classification logic of the audit forces a scope limitation to be recorded when verification fails; the conversion of an unresolved assertion into an accepted one can occur silently in working papers, judgement calls, and system defaults. Audit theory names the epistemic condition but supplies no formal classification architecture that makes the condition impossible to ignore.

\subsection{Accounting Information Systems and the Automation of Audit Classification}
\label{sec:lit-ais}

The accounting information systems literature relocates the verification problem from human judgement to system architecture. Continuous auditing and continuous monitoring research, initiated by Vasarhelyi and Halper's work on continuous process auditing and developed extensively since, envisions audit assurance produced by automated rules operating over transaction streams \citep{vasarhelyi1991, kogan1999, alles2006}. In such systems, classification is performed by exception logic: transactions that violate configured rules are flagged, and flagged items enter queues for investigation.

This literature has documented, though not always in these terms, the conditions under which automated exception handling silently converts verification failure into acceptance. Exception volumes routinely exceed investigative capacity, and prioritisation schemes, thresholds, and alert-fatigue mitigation determine which flags receive attention \citep{alles2006, issa2016}. This exception-prioritisation problem remains an active thresholding question rather than a settled one: recent work formalises how a continuous auditing system should set the threshold that separates flags routed to human review from those cleared without it, treating threshold choice itself as a decision with quantifiable cost \citep{svanberg2025}, which sharpens rather than resolves the question this paper asks -- what state a system should record for the transactions any threshold necessarily excludes from review. More recent work on audit data analytics shows that auditors, managers, and regulators jointly shape how analytic tools are adopted and relied upon, so that the classification problem now extends across the full engagement rather than residing in any single monitoring system \citep{austin2021}. An exception that is deprioritised remains epistemically unverified. Yet because operations continue and dashboards are cleared, the surrounding organisation treats the underlying assertion as accepted. The dashboard state, cleared, is read as the epistemic state, verified, when in fact the two have come apart. This is silent conversion implemented in workflow rather than in judgement.

The AIS literature thus sees the mechanism at ordinary scale. What it lacks is a classification vocabulary in which deprioritised but uninvestigated exceptions occupy a distinct, durable, consequence-bearing state rather than disappearing into workflow completion.

\subsection{Missing Data, SQL Unknown, and the Loss of Uncertainty at the Application Layer}
\label{sec:lit-sql}

At the data layer, accounting information systems already possess a primitive version of the machinery this paper proposes. SQL's treatment of missing information through \textsf{NULL} implements a three-valued logic in which comparisons involving missing data evaluate to \textsf{Unknown} rather than to true or false \citep{codd1979}. The database layer, in other words, refuses to convert absence of information into an affirmative truth value.

The problem arises above the database. Application layers, reporting middleware, dashboards, and audit workflows routinely collapse \textsf{Unknown} into a benign binary output: a null risk score rendered as zero, a missing confirmation aggregated as no exception, a filter predicate that silently excludes unresolved rows so that a completeness check passes. The uncertainty that the data layer carefully preserved is destroyed in transit to the point of audit consequence. This is the entire thesis of the present paper in miniature. The technical capacity to represent verification failure exists and is decades old; what is missing is the architectural discipline to preserve that state through every layer that stands between evidence and audit conclusion.

\subsection{Runtime Verification and Multi-Valued Verdicts}
\label{sec:lit-rv}

Formal methods research confronted a structurally identical problem in runtime verification. A monitor observing a finite execution trace must often render a verdict about a property, such as a temporal logic formula, whose truth may depend on unobserved future behaviour. Classical two-valued semantics forces a premature commitment. The response was to enrich the verdict domain. The logic LTL$_3$ assigns to each finite trace one of three values, true, false, or inconclusive, where the inconclusive verdict records precisely that the observations made so far do not determine the property \citep{bauer2011}. Refinements such as RV-LTL further subdivide the inconclusive region into presumably true and presumably false, permitting operational guidance without permitting the presumptive values to masquerade as definitive ones \citep{bauer2010, leucker2009}.

Runtime verification therefore supplies exactly the formal precedent the audit problem requires: a principled refusal to assign premature truth values on incomplete observations, a machine-checkable semantics for the inconclusive state, and a demonstration that multi-valued verdicts are implementable in monitors that run in production systems. What runtime verification has not done is apply this machinery to audit evidence, where the incomplete trace is not a program execution but an evidence record, and where the consequence of a premature verdict is an audit opinion rather than a property violation report.

\subsection{The Research Gap}
\label{sec:lit-gap}

The literatures reviewed above partition the problem among themselves. Audit theory recognises evidential insufficiency and provides the modified opinion as an output channel, but supplies no formal classification architecture that makes silent conversion structurally impossible. Accounting information systems research studies automated monitoring and exception processing, and documents the workflow conditions under which unverified items are operationally absorbed as accepted, but lacks a classification vocabulary with durable indeterminate states. Database systems preserve missing information as \textsf{Unknown} at the data layer but permit application layers to destroy that state before it reaches audit consequence. Runtime verification provides multi-valued, machine-checkable semantics for inconclusive verdicts on incomplete observations but has not been applied to audit evidence.

What is missing is an accounting information systems methodology that joins these resources into a default-free audit classification architecture: one in which verification boundaries are explicit, indeterminate states are first-class and durable, residual risk is assigned rather than orphaned, and default-freedom is machine-checkable. The gap can be stated in one sentence: to our knowledge, none of these literatures gives accounting information systems a classification architecture that preserves \textsf{Undefined} from the point of verification failure through to audit consequence. Boundary Discipline, developed in the remainder of this paper, is proposed as that architecture. Table~\ref{tab:litmap} summarises the mapping.

\begin{table}[htbp]
\centering
\small
\begin{tabular}{>{\raggedright\arraybackslash}p{2.6cm}>{\raggedright\arraybackslash}p{3.4cm}>{\raggedright\arraybackslash}p{3.6cm}>{\raggedright\arraybackslash}p{4.2cm}}
\toprule
\textbf{Domain} & \textbf{Handles uncertainty as} & \textbf{Failure mode} & \textbf{Boundary Discipline adds} \\
\midrule
Audit evidence & Scope limitation and modified opinion & Judgement may absorb insufficiency & Explicit \textsf{Undefined} and residual ownership \\
\addlinespace
AIS and continuous auditing & Exception queues and thresholds & Ignored exceptions become operationally accepted & Durable unresolved state \\
\addlinespace
SQL and database logic & Three-valued \textsf{Unknown} & Application layer collapses \textsf{Unknown} to benign output & Preservation through audit consequence \\
\addlinespace
Runtime verification & Inconclusive finite-trace verdict & Not translated into audit systems & Multi-valued audit classification \\
\addlinespace
Formal methods & Machine-checked specifications & Rarely tied to audit evidence & Machine-checkable default-freedom \\
\bottomrule
\end{tabular}
\caption{How each domain handles uncertainty, where it fails, and what Boundary Discipline adds.}
\label{tab:litmap}
\end{table}

\section{Boundary Discipline: Core Concepts}
\label{sec:concepts}

This section defines the conceptual apparatus. The two-level vocabulary introduced informally above is now made precise, followed by the accounting instruments that make the vocabulary operational.

\begin{definition}[Verification Failure]
\label{def:vfail}
Verification failure occurs when a system cannot determine whether a claim is verified or refuted within the evidence, authority, and observation boundaries available to it.
\end{definition}

\begin{definition}[Silent Conversion]
\label{def:silent}
Silent conversion occurs when a verification failure is operationally treated as acceptance without an explicit \textsf{Undefined}, modified, blocked, or escalated classification.
\end{definition}

\begin{definition}[Implicit Default]
\label{def:default}
An implicit default is a classification rule that applies when explicit verification procedures fail to produce a determinate result, characterised by three properties: it is not formally specified in the verification protocol; it produces affirmative rather than indeterminate outputs; and its application is not recorded in the verification audit trail.
\end{definition}

Silent conversion is the observable failure mode; the implicit default is the local mechanism by which it occurs. The distinction matters because remediation targets differ. Silent conversion can be produced by an explicit but unrecorded judgement call, by an exception queue whose deprioritisation logic no one reviews, by a middleware coercion of \textsf{NULL} to zero, or by a formally specified rule whose affirmative output is simply wrong for the input class. Boundary Discipline aims at the failure mode, not merely at any one mechanism.

\subsection{The Verification Balance Sheet}
\label{sec:vbs}

Every verification regime can be characterised by a balance sheet that accounts for what verification actually accomplishes. In this architecture the balance sheet is not a formal-kernel data structure; it is an accounting instrument, produced and reviewed by named organisational roles, with its own creation point, review point, and version identifier, exactly as a working paper is.

\begin{definition}[Verification Balance Sheet]
\label{def:vbs}
A Verification Balance Sheet consists of five components. The \emph{Fiat Ground} comprises the assumptions the verification system takes as given without further verification. The \emph{Verified Target} is a precise statement of what successful verification establishes and nothing more. The \emph{Procedure Cost} comprises the resources, in time, money, computation, and institutional overhead, required to complete verification. The \emph{Residuals} comprise what remains unverified after successful completion, including both known limitations and structural blind spots. The \emph{Trust Chain} is the delegation structure showing how the procedure connects to the fiat ground.
\end{definition}

The balance sheet enforces honesty about verification scope. The Verified Target must be stated with sufficient precision that it is clear what has and has not been established, and the Residuals must enumerate what falls outside the verification boundary, especially items that users might mistakenly assume are covered. Table~\ref{tab:vbsfields} states this as a governance record rather than only a definition: what each component requires as a matter of engagement documentation, who is responsible for it, and whether it is machine-generated from a boundary specification's evaluation, human-supplied as professional judgement, or both.

\begin{table}[htbp]
\centering
\small
\begin{tabular}{>{\raggedright\arraybackslash}p{1.8cm}>{\raggedright\arraybackslash}p{2.6cm}>{\raggedright\arraybackslash}p{2.6cm}>{\raggedright\arraybackslash}p{2.6cm}>{\raggedright\arraybackslash}p{2.2cm}}
\toprule
\textbf{Component} & \textbf{Required field} & \textbf{Responsible role} & \textbf{Creation / review point} & \textbf{Source} \\
\midrule
Fiat Ground & assumptions taken as given, with authority relied upon & engagement partner or manager & set at procedure design; reviewed at planning & human-supplied \\
\addlinespace
Verified Target & the exact proposition the procedure, if satisfied, establishes & engagement manager, per procedure & set at procedure design; fixed once the boundary specification is versioned & human-supplied, then machine-enforced by the boundary specification \\
\addlinespace
Procedure Cost & time, money, computation, institutional overhead consumed & engagement manager & recorded at procedure completion & human-supplied \\
\addlinespace
Residuals & what remains unverified after completion, including structural blind spots & preparer, reviewed by manager & recorded at procedure completion; every entry also written to the Risk Assignment Register (Section~\ref{sec:rar}) & mixed: the boundary evaluation records \emph{which} preconditions failed; the preparer records \emph{what that implies is still unestablished} \\
\addlinespace
Trust Chain & the delegation path from the procedure's output back to the Fiat Ground & engagement manager, reviewed by quality control & set at procedure design; revisited if any link is later found compromised & human-supplied \\
\bottomrule
\end{tabular}
\caption{The Verification Balance Sheet as an engagement record: required fields, responsible role, when each is created or reviewed, and whether it is machine-generated, human-supplied, or both. The version identifier binding a specific balance sheet to a specific boundary specification version is carried by the boundary specification itself (Section~\ref{sec:yaml}), not by a separate field here.}
\label{tab:vbsfields}
\end{table}

The relationship to the evidence record and to the audit conclusion is direct and is what distinguishes this instrument from a working paper narrative: the Verified Target is the proposition the boundary specification's postcondition licenses (Section~\ref{sec:preconditions}), the Residuals are what the same specification's failed preconditions name, and the audit conclusion may rely affirmatively on the Verified Target alone -- never, by the Central Invariant (Principle~\ref{prin:invariant}), on anything the Residuals line discloses.

A passport check illustrates the discipline outside the audit domain, where the components are easier to see without the weight of the running case. Its Fiat Ground includes the authority of the issuing government and the integrity of the issuance process. Its Verified Target is narrow: a government issued a document with this name and photograph to someone, and the presenter resembles the photograph. Its Residuals include whether the original issuance was correct, whether the document has been tampered with in undetectable ways, and whether the recorded name corresponds to any fact about identity beyond what the government chose to record.

Section~\ref{sec:wirecard} gives a worked Verification Balance Sheet for the Wirecard cash-existence assertion itself, populated against the fields of Table~\ref{tab:vbsfields} rather than only narrated in prose.

\subsection{The Risk Assignment Register}
\label{sec:rar}

Verification residuals represent risk, and Boundary Discipline requires that every residual either be explicitly assigned to an identified owner or be flagged, visibly and durably, as orphaned -- the owner field is optional, not the escalation. This register is where the architecture's exclusive AIS contribution is clearest: the companion formal development validates that residuals are emitted and preserved; it does not, and does not claim to, perform organisational risk assignment, which is what this section defines.

\begin{definition}[Risk Assignment Register]
\label{def:rar}
A Risk Assignment Register associates each verification residual with answers to three questions: what is the risk, who owns it, and what is their capacity to bear it. If no owner can be identified for a residual, that fact is itself a first-class finding requiring escalation.
\end{definition}

Table~\ref{tab:rarschema} states this as a record schema an AIS implementation can actually build. The columns fall into two groups: those a formal classification kernel can populate mechanically from a boundary evaluation, and those an organisational governance layer must populate under its own procedures.

\begin{table}[htbp]
\centering
\small
\begin{tabular}{>{\raggedright\arraybackslash}p{3.6cm}>{\raggedright\arraybackslash}p{5.4cm}>{\raggedright\arraybackslash}p{2.6cm}}
\toprule
\textbf{Field} & \textbf{Description} & \textbf{Populated by} \\
\midrule
Residual identifier & stable, unique key for this register entry & kernel \\
Assertion identifier & the material assertion the residual concerns & kernel \\
Engagement identifier & the audit engagement this residual belongs to & AIS layer \\
Boundary-specification identifier and version & which declared boundary, and which version of it, produced the classification & kernel \\
Evidence-context identifier & which evidence context the classification was evaluated against & kernel \\
Classification state & the resulting state (Section~\ref{sec:multivalued}) & kernel \\
Failure reason & why the entry was written: not verified, verified but unadmitted, or admission invalid & kernel \\
Materiality & the amount or qualitative basis on which the assertion was judged material & AIS layer (recorded, not derived) \\
Financial-statement area & the account, class of transaction, or disclosure the assertion supports & AIS layer \\
Proposed owner & the individual or function nominated to own the residual risk & AIS layer \\
Owner acceptance status & whether the proposed owner has accepted ownership & AIS layer \\
Owner capacity & the accepting owner's assessed capacity to bear the risk & AIS layer \\
Orphan status & true whenever no owner has both been proposed and accepted & kernel detects and reports; AIS layer resolves \\
Escalation authority & the role or committee an orphaned or unaccepted residual is escalated to & AIS layer \\
Disclosure status & whether, and where, the residual has been disclosed externally & AIS layer \\
Creation timestamp & when the entry was written & kernel \\
Review timestamp & when the entry was last reviewed by the AIS layer & AIS layer \\
Closure reference & a future governance layer's record of resolution, if any (Section~\ref{sec:limitations}: not implemented here) & AIS layer, future work \\
Historical provenance & the full append-only history of the entry, including every review that left it open & kernel (append-only log) and AIS layer (review annotations) \\
\bottomrule
\end{tabular}
\caption{Risk Assignment Register schema. ``Kernel'' denotes what the companion formal development mechanically computes or preserves; ``AIS layer'' denotes what only an organisational governance process, outside the kernel, can supply.}
\label{tab:rarschema}
\end{table}

The kernel boundary in Table~\ref{tab:rarschema} is stated once here and holds throughout the paper: the kernel emits the residual and may represent an optional owner field and query for orphaned entries; it does not assign an owner, does not determine owner capacity, and does not close a residual. The organisational AIS layer performs those functions, under governed procedures this paper specifies architecturally (Section~\ref{sec:roles}) but that the companion kernel neither implements nor claims to.

Table~\ref{tab:rarentries} gives three illustrative entries, populated against the schema above. They are illustrative rather than data from any real engagement, and are labelled as such throughout: no entry below reports an amount, a name, or an outcome drawn from an actual audit file.

\begin{table}[htbp]
\centering
\small
\begin{tabular}{>{\raggedright\arraybackslash}p{2.2cm}>{\raggedright\arraybackslash}p{3.4cm}>{\raggedright\arraybackslash}p{3.4cm}>{\raggedright\arraybackslash}p{3.4cm}}
\toprule
\textbf{Field} & \textbf{1. Wirecard cash existence (illustrative)} & \textbf{2. Deprioritised CA exception (illustrative)} & \textbf{3. Orphaned residual (illustrative)} \\
\midrule
Assertion & EUR 1.9bn Philippine trustee balance exists & Transaction \#5 violates no defined control & Cross-jurisdiction related-party balance exists \\
Classification state & Undefined & Undefined & Undefined \\
Failure reason & not verified: route not auditor-controlled & not verified: no review completed before threshold cut-off & not verified: no declared procedure covers this assertion class \\
Materiality & material to cash and cash equivalents & below individual materiality; material in aggregate with seven similar entries & material \\
Proposed owner & Group audit partner & Engagement senior, continuous-auditing queue & \emph{(none proposed)} \\
Owner acceptance status & not recorded (illustrative gap) & accepted & N/A \\
Orphan status & no (owner proposed, acceptance not yet recorded) & no & \textbf{yes} \\
Escalation authority & audit committee & N/A at this materiality & audit committee, immediately (per Principle~\ref{prin:ownership}) \\
Disclosure status & not disclosed (illustrative) & internal only & escalated, not yet disclosed \\
\bottomrule
\end{tabular}
\caption{Illustrative Risk Assignment Register entries. Entry 1 corresponds to Section~\ref{sec:wirecard}'s worked case; entry 2 corresponds to one of the eight deprioritised transactions in Section~\ref{sec:ca-case}; entry 3 illustrates Principle~\ref{prin:ownership}'s orphan-escalation branch, which entries 1 and 2 do not exercise.}
\label{tab:rarentries}
\end{table}

A residual entry is emitted with no owner: the classification mechanism that writes it does not itself assign one. Ownership is an answer supplied afterward, by whatever governance process this discipline is embedded in; what the register (and, in the companion formal development, the kernel's orphan query) does mechanically is \emph{find} residuals that still have no owner and make that fact visible, not \emph{assign} an owner to them. The distinction matters for what Appendix~\ref{app:requirements} and Section~\ref{sec:implications} can honestly claim: this paper's discipline makes orphaned risk detectable and durable, not automatically resolved.

\begin{principle}[Risk Ownership]
\label{prin:ownership}
In any verification regime, every residual must be either assigned to an identified risk owner with capacity to bear it, or flagged as orphaned and thereby escalated and disclosed, or explicitly priced by parties who choose to trade it.
\end{principle}

Orphaned risk is systemic fragility waiting to materialise. In the Wirecard case, examined in Section~\ref{sec:wirecard}, the risk that \officialeuro 1.9 billion did not exist was a verification residual that no party owned: the auditor treated it as outside scope, the home regulator treated the entity as outside its banking mandate, the foreign banks had no relationship to protect, and investors assumed the audit provided coverage. The risk accumulated in the gap between verification regimes until the fraud was revealed.

Once residuals are explicit, they can also be priced, insured, escalated, or disclosed. The economic treatment of verification value, in which the value created by a procedure is the measurable difference between pre-verification and post-verification risk exposure, is developed briefly in Section~\ref{sec:implications} as an implication of the framework rather than as a core component.

\section{Design Principles for Default-Free Audit Classification}
\label{sec:principles}

The following principles govern the architecture developed in the next section. They are stated as constraints that any default-free audit classification system must satisfy.

\begin{enumerate}
\item \textbf{No affirmative classification outside a verified boundary.} The classifier may, and in this architecture does, return a determinate output for every input; what it must not do is return an \emph{affirmative} one for an input over which it has no verified authority.
\item \textbf{Unverified is not accepted.} Failure to verify must return a distinct classification state, never an affirmative one.
\item \textbf{Presumptive states cannot bind audit consequence.} A classification of likely acceptable may guide triage and workflow, but it cannot support an unqualified opinion or any equivalent binding conclusion.
\item \textbf{Every residual must be owned, disclosed, or escalated.} No verification residual may remain unassigned without that fact itself being recorded and escalated.
\item \textbf{Admission requires a distinct verifier identity.} The registered identity proposing acceptance of a claim cannot also be the identity that admits it; admission requires a distinct, registered identity holding verifier authority.
\item \textbf{The audit trail must record classification failure.} Missing evidence, weak evidence, jurisdictional obstruction, and unsupported management representation must leave durable traces that survive workflow completion.
\end{enumerate}

An implicit default in the sense of Definition~\ref{def:default} necessarily violates the first, second, and sixth principles directly -- it produces an affirmative output outside verified authority, that output is accepted rather than distinct, and its application is unrecorded -- and it need not, on its own, violate the third, fourth, or fifth; those are separate failure modes a system can exhibit independently of whether it contains an implicit default at all. Each principle is individually motivated by the failure modes documented in Section~\ref{sec:lit}: the first two by the audit evidence literature, the third by the runtime verification distinction between presumptive and definitive verdicts, the fourth by the orphaned-risk pattern, the fifth by the structural requirement that admission be separated from proposal, and the sixth by the observation that silent conversion thrives precisely where classification failure leaves no trace.

\section{A Default-Free AIS Classification Architecture}
\label{sec:architecture}

Boundary Discipline can now be described as an accounting information systems architecture: an information pipeline in which classification is explicit, multi-valued, and admission-gated.

\subsection{The Pipeline}
\label{sec:pipeline}

A claim enters the system as a \emph{claim packet}: the assertion itself together with its evidence packet, comprising the documents, confirmations, log extracts, or automated check results offered in its support, and the metadata needed for boundary evaluation, including assertion class, jurisdiction, materiality, and evidence grades. A \emph{boundary specification}, maintained as versioned machine-readable data, determines whether the attached evidence is admissible for that claim type: which verification methods produce valid evidence for which assertion classes, in which jurisdictions, at which materiality thresholds. The \emph{classification function} evaluates the claim packet against the boundary specification and returns one of the states defined in Section~\ref{sec:multivalued}; as Section~\ref{sec:multivalued} makes precise, the generic, mechanised classification function returns \textsf{Verified} or \textsf{Undefined}, and the remaining four states are vocabulary the surrounding architecture accepts and propagates, not values that function derives on its own. An independent \emph{admission verifier}, distinct from whatever process generated or proposed the claim, determines whether the classification is admitted for audit consequence. All non-affirmative outcomes write entries to the Risk Assignment Register, and only admitted \textsf{Verified} claims may bind the audit conclusion.

The architectural point, borrowed from fail-closed admission control in verified systems, is the strict separation of proposal from admission. Execution and proposal may be performed by any process, human or automated, including generative tools. Admission is performed only by the verifier, only against the declared boundary specification, and only into the declared classification states. A claim packet that fails admission is not rejected into silence; it is classified, registered, and escalated according to its state.

One qualification on the admission verifier applies throughout and is stated here rather than left implicit. In the companion formal development, the verifier and proposer are each identified by an identifier the caller supplies and registers, and a certificate is checked against those identifiers -- it names the proposal it admits, and the registry is consulted, by identifier, for role and authority. That development does not itself establish that the identifiers a caller supplies are unique. Two proposals sharing an identifier make a certificate drawn up for one also validate against the other; two registry entries sharing a process identifier make role and authority lookups depend on which entry is found first rather than on a well-defined population. Uniqueness of these identifiers is accordingly a caller responsibility this architecture assumes rather than enforces, in the same sense the boundary specification and materiality threshold are caller-supplied without a well-formedness check.

\subsection{Organisational Responsibility Model}
\label{sec:roles}

The pipeline above names two roles, proposer and admitting verifier, at the level a formal kernel can check: distinct registered identities. An AIS implementation needs more roles than that to actually govern an engagement, and this section names them, distinguishes what each is responsible for, and states plainly what registry-level identity separation does and does not establish about them.

\begin{itemize}
\item \textbf{Evidence provider.} Supplies the underlying documents, confirmations, log extracts, or automated check results a claim packet's evidence context is built from. May be internal (an AIS monitor) or external (a bank, a counterparty).
\item \textbf{Claim proposer.} Constructs a claim packet and proposes a classification for admission. Corresponds to the kernel's registered \texttt{Proposer} role.
\item \textbf{Classification service.} Evaluates a claim packet against the declared boundary specification and returns \textsf{Verified} or \textsf{Undefined}; this is the one role a formal kernel actually performs, mechanically, without organisational discretion.
\item \textbf{Registered admission verifier.} Reviews a proposed classification and either admits it, binding it to audit consequence, or does not. Corresponds to the kernel's registered \texttt{Verifier} role.
\item \textbf{Residual owner.} Accepts responsibility for a specific residual's risk, once proposed by governance and recorded in the Risk Assignment Register (Section~\ref{sec:rar}).
\item \textbf{Owner-capacity reviewer.} Independently assesses whether a proposed owner actually has the authority, resources, or standing to bear the risk they are being asked to own -- a check the register records but does not itself perform.
\item \textbf{Escalation authority.} Receives orphaned residuals and residuals whose owner-capacity review failed; typically the audit committee or an equivalent governance body for material items (Table~\ref{tab:rarentries}, entry 3).
\item \textbf{Audit-committee or governance recipient.} Receives the aggregate Risk Assignment Register and Verification Balance Sheet summaries this architecture's reporting implications describe (Section~\ref{sec:implications}), distinct from the escalation authority for any single item.
\item \textbf{External inspection authority.} Consumes the boundary coverage, undefined-region materiality, and residual assignment measures of Section~\ref{sec:inspection} from outside the engagement, without access to the underlying working papers.
\end{itemize}

These roles may be fulfilled by different people, different systems, or different organisations, and in a small engagement several may be fulfilled by the same person wearing different hats at different times -- that is an organisational design choice this architecture does not make on a firm's behalf. What the architecture does require, and what the admission-separation property Table~\ref{tab:validation} reports (Section~\ref{sec:coq}) formally validates, is narrower: the claim proposer and the registered admission verifier for the \emph{same} classification cannot be the same registered identity. Registry-level identity separation is exactly that and no more. It does not establish, and this paper does not claim, that two distinct registered identities are controlled by different people, different reporting lines, or different organisations; a firm that registers two identities under its own control has satisfied the kernel's check while satisfying none of the organisational independence an audit committee would actually want. Imposing and inspecting that stronger requirement -- who may hold the verifier role, what reporting-line separation it requires, whether rotation or external review is needed -- is the AIS governance layer's job, not the kernel's, and this paper's contribution is naming that job precisely enough to be designed and inspected, not performing it by assumption.

\subsection{Multi-Valued Classification}
\label{sec:multivalued}

The theorem-level core of the framework, presented in Section~\ref{sec:formal}, requires only the binary distinction between \textsf{Verified} and \textsf{Undefined} for affirmative purposes, and this is also the only distinction the generic, mechanised classification function in the companion Coq/Rocq development actually computes: given a declared boundary specification and an evidence context, it returns \textsf{Verified} when some declared procedure's preconditions are satisfied and \textsf{Undefined} otherwise, with no third branch for it to take. Operationally, however, audit workflow benefits from a richer verdict domain, and the runtime verification literature shows how to enrich it without corrupting it. The framework accordingly declares six states below, and the composition and admission machinery of this paper accepts a claim already classified into any of them; but only the two just described are produced by a generic, proved-correct procedure. \textsf{Refuted} and the two presumptive states are legitimate inputs to that downstream machinery, not outputs the classification function derives on its own; no generic, mechanised procedure for \emph{assigning} one of those three from a boundary specification and an evidence context is claimed here.

\textsf{Verified}: the claim falls within a declared boundary and the required evidence grade is present. This is the only state that can support an unqualified audit conclusion.

\textsf{Refuted}: admissible evidence establishes that the claim is false.

\textsf{Undefined}: the claim falls outside all declared verification boundaries, or the required evidence grade is absent, and no admissible procedure can currently resolve it.

\begin{remark}
\label{rem:success}
\textsf{Undefined} is a success state. A classifier that returns \textsf{Undefined} for an input outside its verified domain has computed the correct answer. The failure states of an audit classification system are not its \textsf{Undefined} outputs; they are its affirmative outputs on inputs it had no authority to classify.
\end{remark}

\textsf{Presumptively Verified}: available evidence favours the claim but does not meet the declared grade for the claim's materiality. This state may guide workflow prioritisation. It cannot bind audit consequence.

\textsf{Presumptively Refuted}: available evidence disfavours the claim without admissibly establishing falsity. The same consequence restriction applies.

\textsf{Escalation Required}: the claim exposes a defect in the boundary specification itself, such as an overlap of conflicting requirements or a claim class for which no boundary was declared, and must be raised to the level at which specifications are governed. This is a distinct concern from an orphaned residual (Section~\ref{sec:rar}): an orphaned residual is a classification, of any of the six states above, that additionally lacks an identified risk owner, discovered by a query over the register after the fact, not a defect in the boundary specification that the classification function itself detects or reclassifies as \textsf{Escalation Required}. No generic, mechanised procedure for assigning \textsf{Escalation Required} is claimed here either.

The presumptive states are the audit analogue of the refined inconclusive verdicts in RV-LTL: they encode operational lean without pretending to definitive resolution, and the consequence restriction is what prevents the enrichment from reintroducing silent conversion through the back door. A dashboard may sort its queue by presumptive state; an opinion may not rest on one.

\subsection{Residual Registers and Durable Traces}
\label{sec:registers}

Every classification other than an admitted \textsf{Verified} generates a residual entry recording the claim, the state assigned, the boundary evaluation that produced it, and the risk assignment required by Principle~\ref{prin:ownership}. The register is append-only with respect to workflow: clearing a dashboard, closing a queue item, or completing an engagement phase does not delete or downgrade a residual entry. This single property, durability of the non-affirmative trace across workflow completion, is what distinguishes the architecture from a conventional exception system, and it is the direct structural answer to the absorption pattern documented in Section~\ref{sec:lit-ais}.

\subsection{Composition}
\label{sec:composition}

Audits are made of many procedures, and an opinion depends on many assertions. The architecture is worthless unless correct local classifications compose into a correct conclusion, so the composition rule is stated explicitly.

\begin{principle}[Composition]
\label{prin:composition}
An audit opinion is admissible as unqualified only if every material assertion on which it depends is an admitted \textsf{Verified} classification. Any material dependency that is not both verified and validly admitted blocks the opinion, regardless of how many other dependencies are.
\end{principle}

\begin{lstlisting}[numbers=none]
if any material assertion in deps(Opinion) is Undefined and unresolved
then Opinion is inadmissible as unqualified
\end{lstlisting}

The direction matters, informally: strength does not propagate upward from a majority of well-verified assertions; a single unresolved one blocks admissibility regardless of how many others are Verified. This is not pessimism. It is the observation that an unqualified opinion asserts a conjunction, and a Boolean conjunction is no stronger than its weakest (i.e.\ false) conjunct -- a claim about the two-valued admissible/not-admissible status of each assertion for this purpose, not a claim that the six declared classification states (Section~\ref{sec:multivalued}) form any further order among themselves; no such order is defined or needed here. A methodology that averages evidence quality across assertions, or that permits a preponderance of \textsf{Verified} classifications to carry an \textsf{Undefined} one across the admissibility line, has reintroduced the implicit default at the composition layer after eliminating it at the classification layer.

The dependency set this rule quantifies over is not this architecture's claim about what an audit opinion asserts at the level of individual transactions. An opinion ordinarily provides reasonable assurance over the financial statements as a whole, not a conjunction of every transaction separately inspected. The dependency set is the set of assertions the engagement's own risk and materiality assessment has already scoped as material to that opinion -- an upstream, auditor-supplied determination this architecture takes as given (Section~\ref{sec:limitations}), not one it derives.

\section{Formal Validation of Selected Architectural Requirements}
\label{sec:formal}

The methodology becomes enforceable only when its components are specified precisely enough to admit machine verification. The formalisation target is deliberately narrow and structural: the classification logic that maps evidence types, jurisdictions, materiality thresholds, and control relationships into audit classifications. This is the locus where implicit defaults typically enter. Substantive sampling strategy, statistical sufficiency, valuation modelling, fraud detection, and management intent are not formalised; they remain domains for professional judgement and substantive audit work. Boundary Discipline constrains how such work is allowed to propagate into audit conclusions by forbidding silent defaults when verification is structurally unavailable.

\subsection{Boundary Specification Language}
\label{sec:yaml}

Verification boundaries are specified in a structured, versioned, machine-readable record supporting audit trails and automated validation. Table~\ref{tab:boundaryspec} gives an illustrative specification for the cash-existence assertion class underlying the Wirecard case discussed in Section~\ref{sec:wirecard}; the schema is illustrative and would require domain-specific validation by auditors, regulators, and AIS specialists before production use. Each row binds a verification method to the evidence quality it can establish, the jurisdictions in which it is recognised, the materiality threshold up to which it applies, and whether it is sufficient on its own to support an unqualified opinion.

\begin{table}[htbp]
\centering
\small
\begin{tabular}{>{\raggedright\arraybackslash}p{2.6cm}>{\raggedright\arraybackslash}p{2.6cm}>{\raggedright\arraybackslash}p{1.6cm}>{\raggedright\arraybackslash}p{3.2cm}>{\raggedright\arraybackslash}p{2.6cm}}
\toprule
\textbf{Verification method} & \textbf{Evidence quality} & \textbf{Required for opinion} & \textbf{Applicable jurisdictions} & \textbf{Materiality threshold} \\
\midrule
Direct bank confirmation & Primary & Yes & Germany, United Kingdom, United States, Japan, Australia & Any amount \\
\addlinespace
Third-party trustee certification & Secondary & No (escalation required) & Singapore, Ireland & Below 5\% of total assets \\
\addlinespace
Management representation only & Tertiary & No (opinion modification required) & None declared & Non-material only \\
\bottomrule
\end{tabular}
\caption{Illustrative verification boundary specification for the cash-existence assertion class. Material assets in a jurisdiction not listed against a boundary method that is required for opinion fall in an undefined region: no unqualified opinion may issue, and audit committee and board notification is required.}
\label{tab:boundaryspec}
\end{table}

The specification makes explicit what current practice leaves implicit. When a material balance sits in trustee accounts in a jurisdiction that appears in no boundary's jurisdiction list for material amounts, the classification function cannot return \textsf{Verified}. This is not a failure of the system; it is the system working correctly by refusing to classify outside its verified domain, exactly as Remark~\ref{rem:success} requires. In practice such a specification would be maintained as versioned, machine-readable data (for example, as structured records analogous to the boundary tables an XBRL taxonomy already uses to constrain which tags and contexts are valid for a given filing), rather than as a static table; Table~\ref{tab:boundaryspec} presents its logical content rather than its storage format.

\subsection{Preconditions and Postconditions of Verification Procedures}
\label{sec:preconditions}

The boundary specification becomes mechanical by giving every verification procedure a proof obligation of the same shape as a program statement: a precondition stating what must be true before the procedure is entitled to produce a verification result, and a postcondition stating exactly which classification the procedure is entitled to produce. The question to ask of every audit procedure is the weakest-precondition question: what must hold for this procedure's output to mean what it claims to mean?

\begin{lstlisting}[numbers=none, caption={Precondition and postcondition of direct bank confirmation.}, label={lst:precond}]
Procedure: Direct bank confirmation

Preconditions:
  bank exists as independent counterparty
  confirmation route is controlled by the auditor
  jurisdiction permits direct confirmation
  account identifier is supplied
  response is received through an authenticated channel

Postcondition:
  the cash-existence assertion may be classified Verified
\end{lstlisting}

If any precondition fails, the procedure cannot produce \textsf{Verified}. Not should not: cannot, within the declared system. A confirmation received through a route the auditor does not control, or from a jurisdiction that obstructs direct confirmation, does not weakly satisfy the precondition; it fails it, and the postcondition is unavailable. Silent conversion, in this vocabulary, is a postcondition violation: a procedure whose preconditions failed was nonetheless treated as having discharged its postcondition. An assertion for which no declared procedure's preconditions hold is classified \textsf{Undefined} by construction.

\subsection{Selected Requirements and Their Formal Validation}
\label{sec:theorems}

Four architectural requirements from Sections~\ref{sec:principles}--\ref{sec:composition} admit formal validation without the audit substance around them being formalised: that classification without a satisfied declared procedure cannot be affirmative, that admission requires a distinct registered verifier identity, that a material failed dependency durably emits a residual, and that an opinion cannot be admitted unqualified while any material dependency remains unresolved. Table~\ref{tab:validation} states each requirement, the companion kernel property that validates it, and what that validation means architecturally.

\begin{table}[htbp]
\centering
\small
\begin{tabular}{>{\raggedright\arraybackslash}p{2.6cm}>{\raggedright\arraybackslash}p{4.8cm}>{\raggedright\arraybackslash}p{4.6cm}}
\toprule
\textbf{AIS requirement} & \textbf{Supporting kernel property} & \textbf{Architectural significance} \\
\midrule
No affirmative default & Verified output has a satisfied declared procedure & Unresolved evidence remains visible \\
\addlinespace
Admission separation & Certificate requires a distinct registered verifier identity & Proposal cannot admit itself inside the declared registry \\
\addlinespace
Residual preservation & Material failed dependencies emit append-only entries & Workflow completion cannot erase the historical trace \\
\addlinespace
Opinion restriction & \textsf{Unqualified} requires every material dependency to pass & Partial evidence cannot be averaged into acceptance \\
\bottomrule
\end{tabular}
\caption{Selected architectural requirements and the companion kernel property that formally validates each; Section~\ref{sec:coq} states the scope and limits of this validation.}
\label{tab:validation}
\end{table}

These results validate the implementation of the selected architectural requirements. They do not validate the normative correctness of the requirements -- that this is the right way to define an admissible opinion is an argument this paper makes in prose (Sections~\ref{sec:principles}--\ref{sec:composition}), not something a proof assistant adjudicates -- nor the truth of caller-supplied evidence and materiality judgements; Section~\ref{sec:limitations} states this boundary in full.

Two further requirements from ordinary audit practice are architecturally motivated but not among the properties Table~\ref{tab:validation} reports. Jurisdictional exclusivity requires that no entity fall under multiple regulatory boundaries with conflicting verification requirements without explicit reconciliation, so that gaps between regulatory regimes are flagged rather than silently tolerated. Related-party transitivity requires that control relationships compose transitively, so that functional control exercised through intermediaries cannot be concealed by nominal ownership structures.

\subsection{Formal Validation}
\label{sec:coq}

Formal validation here is not directed at the audit itself, and it does not extend to auditing standards, sampling, or professional judgement. Its target is narrow: whether the classification and admission requirements of Table~\ref{tab:validation} can be implemented, together, by an executable procedure that never violates them -- a question about the artefact's construction, not about which requirements an audit should impose. That narrower question is answered by a companion Coq/Rocq development and its extracted OCaml classifier, cited in the Data availability statement, which the present manuscript treats as supporting evaluation evidence rather than as its own technical contribution. The companion manuscript, \emph{A Discipline of Partial Audit Classification: A Machine-Checked Kernel for Default-Free Audit Decisions}, states the concrete data structures, function definitions, theorem statements, module dependencies, declaration counts, \texttt{Print Assumptions} results, trusted-computing-base analysis, and extraction details this manuscript deliberately does not reproduce.

At this manuscript's level of detail: the companion kernel computes \textsf{Verified} or \textsf{Undefined} from a declared boundary specification and an evidence context, with no third branch for the generic classifier to take. Valid admission requires a certificate naming the correct proposal, presented by a registered identity holding verifier authority, distinct from the registered proposer. A material dependency that is not both verified and validly admitted automatically emits a residual entry that no subsequent workflow operation removes. An opinion decision returns \textsf{Unqualified} if and only if every material dependency passes both checks. The implementation was machine-checked: every stage is proved sound individually, and the end-to-end result composes those proofs rather than being proved independently. Six adversarial fixtures in the companion development exercise the certificate check specifically -- a proposer naming itself as verifier, an unregistered proposer, and a proposer registered with the wrong role -- each rejected by the mechanism Table~\ref{tab:validation}'s admission-separation row reports.

What this validation does not establish bears directly on how this architecture should be read, not only on the companion repository: evidence authenticity, correct materiality selection, organisational independence between distinct registered identities (Section~\ref{sec:roles}), uniqueness of caller-supplied identifiers, residual closure, and the correctness of the extraction and compilation steps that produce the running executable all remain outside the proved boundary, exactly as Section~\ref{sec:limitations} states for this architecture as a whole.

\section{Illustrative Case: Wirecard and Material Cash Verification}
\label{sec:wirecard}

The Wirecard collapse can now be read through the framework. Three boundary violations were operative. The confirmation-route boundary was never enforced: the declared procedure for material cash existence, Table~\ref{tab:boundaryspec}'s direct bank confirmation, requires a route the auditor controls, and the route actually used ran through a third-party trustee the auditor did not control, for years, before direct confirmation was ever attempted -- the Philippine banks were not, in the end, jurisdictionally unreachable; they answered as soon as they were asked directly. The evidence quality boundary between primary bank confirmation and tertiary documentation, such as screenshots and intermediary-provided statements, was never formally enforced against the materiality involved. The control boundary between independent third-party processors and functionally controlled entities was never formally defined, permitting opportunistic classification of Asian partners that derived their entire revenue from Wirecard.

Under Boundary Discipline, the material cash-existence assertion is admissible only if primary confirmation is available within the applicable jurisdiction list of the boundary specification. When that confirmation is unavailable, the preconditions of Section~\ref{sec:preconditions} fail, the no-affirmative-default property (Table~\ref{tab:validation}) forces the classifier into a non-affirmative state, and the composition rule (Principle~\ref{prin:composition}) makes an unqualified opinion formally inadmissible given that state. The claim being made here is deliberately narrow, and conditional. Boundary Discipline would not have prevented the fraud, and it would not have compelled any particular exercise of professional judgement; nor can the kernel establish that a claim packet's recorded facts are themselves true -- a compromised or complicit process could still assert ``route controlled by auditor'' falsely, and no mechanism described here detects that. What the framework offers is narrower: had the claim packet truthfully recorded what the confirmation actually was -- a route not controlled by the auditor, evidence of tertiary rather than primary grade -- the declared preconditions of Table~\ref{tab:boundaryspec} would have failed, and the classifier would have returned \textsf{Undefined} rather than permitting the postcondition to be treated as discharged. That conditional claim is what the classification-theoretic reading of Wirecard establishes: not that the fraud would have been prevented, but that the legacy decision procedure was not default-free even under an honest account of its own evidence, and that a default-free one would have recorded the obstruction rather than let it pass as acceptance.

The Verification Balance Sheet for the bank balance assertion makes the accounting explicit, populated here against the fields of Table~\ref{tab:vbsfields}. The Fiat Ground comprises the bank confirmation channel and the assumption that a route represented as auditor-controlled actually is. The Verified Target is the existence of specific cash balances at named institutions, and nothing more. The Procedure Cost comprises direct confirmation and audit delay. The Residuals comprise trustee independence and the authenticity of the confirmation route -- exactly the residual that, left unassigned for years, is what the fraud exploited. The Trust Chain runs from auditor to trustee to bank to the public audit opinion, and it is the trustee link in that chain, not any jurisdictional one, where the fraud lived: every element the fraud exploited traces to that one corrupted link -- the confirmation route was not what it was represented to be, the trustee was not independent, and the residual this created was owned by no one -- and to that link having gone unassigned in a Risk Assignment Register for years before the banks were ever asked directly. A methodology that had required this balance sheet, and had required each residual to appear in a Risk Assignment Register with an identified owner, would have recorded the evidential obstruction as a first-class classification outcome rather than allowing it to disappear into workflow.

Table~\ref{tab:rarentries}'s first illustrative entry states what a Risk Assignment Register would have carried for this assertion: a residual proposed to the group audit partner as owner, with acceptance status unrecorded -- itself an illustrative gap, since an unaccepted proposed owner is functionally orphaned under Principle~\ref{prin:ownership} and should have triggered audit-committee escalation on exactly the timeline this case shows did not happen. ISA 505 treats external confirmation, sought and controlled by the auditor, as ordinarily the more reliable evidence for an existence assertion of this kind \citep{isa505}; ISA 580 is explicit that management or intermediary representations alone do not provide sufficient appropriate evidence \citep{isa580}; and had the residual been recorded rather than absorbed, the scope limitation it represents would have been the kind of evidential obstruction ISA 705's qualified-opinion or disclaimer channel exists for \citep{isa705}, not an unqualified opinion issued around it. None of this required detecting the fraud in advance. It required the architecture -- a Risk Assignment Register with owner acceptance tracked, not merely proposed, and an escalation authority that acts on an unaccepted material residual -- that this paper's organisational responsibility model (Section~\ref{sec:roles}) specifies and the legacy engagement did not have.

\section{AIS Use Case: Silent Conversion in Continuous Auditing}
\label{sec:ca-case}

The same structure appears at ordinary scale in continuous auditing. Consider a monitoring system whose configured rules flag transactions for investigation, and whose exception volume exceeds investigative capacity, so that a prioritisation scheme determines which flags are pursued. A worked example makes the mechanism concrete. Suppose the system flags 10{,}000 transactions in a reporting period and a thresholding layer routes 200 of them to human review. The remaining 9{,}800 are not verified; they are merely deprioritised. If the dashboard reports that 200 exceptions were reviewed with no material issue found, while the 9{,}800 disappear from consequence, the system has silently converted unreviewed exceptions into accepted background. Exceptions that are deprioritised for practical reasons remain epistemically unverified. Yet because operations continue and the dashboard is cleared, the system, and the organisation reading it, treats the underlying assertions as accepted by default. The dashboard is cleared, but the assertions have not been verified. Silent conversion here occurs not through human judgement on any single item but through thresholding, alert fatigue, and the mechanics of the exception queue.

The mechanism can be written as a three-line algorithm. The undisciplined system runs:

\begin{lstlisting}[numbers=none]
for each exception e:
    if priority(e) < threshold:
        clear(e)
\end{lstlisting}

The disciplined system runs:

\begin{lstlisting}[numbers=none]
for each exception e:
    if priority(e) < threshold:
        # state(e) is already Undefined -- classification against
        # the declared boundary computed this, not this workflow step
        residual_register.add(e)
\end{lstlisting}

with the governing inequality:

\begin{lstlisting}[numbers=none]
Undefined != Accepted
\end{lstlisting}

That single line is the continuous auditing case in full. Prioritisation is legitimate, capacity is finite, and triage is not the error. The error is the disappearance of the deprioritised exceptions from the record of what remains unestablished. A deprioritised exception is not assigned a separate \textsf{Unreviewed} state that this workflow step invents; it remains \textsf{Undefined}, exactly as classification against the declared boundary already computed it (Section~\ref{sec:multivalued}), and the residual entry records that computed state rather than a new one.

The Verification Balance Sheet for a continuous auditing exception reads as follows. The Fiat Ground comprises the transaction logs and the rule configuration. The Verified Target is whether a specific transaction violates a defined control, and nothing more. The Procedure Cost comprises the automated flag plus human investigation. The Residuals comprise deprioritised exceptions, threshold tuning, and false negatives. The Trust Chain runs from the AIS monitor to the audit dashboard to management assurance.

Under Boundary Discipline, a deprioritised exception is classified \textsf{Undefined}: without the review its declared boundary requires, no procedure's preconditions are satisfied, and the generic classifier has no third branch to take (Section~\ref{sec:multivalued}). It generates a durable residual entry that survives the clearing of the dashboard. The aggregate of such entries, visible in the Risk Assignment Register, is itself audit evidence about the assurance the monitoring system actually provides, as distinct from the assurance its cleared dashboard appears to provide. The data layer already knows how to do this, as Section~\ref{sec:lit-sql} observed; the discipline consists in refusing to let any layer between the flag and the opinion collapse the state.

Table~\ref{tab:rarentries}'s second illustrative entry states what one of the 9,800 deprioritised transactions carries in the Risk Assignment Register: owned by the engagement senior responsible for the continuous-auditing queue, accepted, below individual materiality but material in aggregate with the other seven the same worked fixture exercises. Aggregation is precisely where management reporting and escalation thresholds belong and where the classification/workflow distinction matters most: no single deprioritised transaction need reach the audit committee, but the aggregate undefined-region materiality across all 9,800 (Section~\ref{sec:inspection} states how this is measured) is exactly the indicator that should, at a firm-defined threshold, route to the escalation authority (Section~\ref{sec:roles}) regardless of whether any individual item does. A monitoring system that reports "200 reviewed, no material issue found" without also reporting that aggregate has re-committed the silent conversion this paper's central invariant forbids, one dashboard summary at a time.

\section{Implications for Standards, Audit Reports, and Inspection}
\label{sec:implications}

Boundary Discipline could inform a future auditing standard, an inspection protocol, or an internal firm methodology; the framework does not depend on any one of these adoption paths. This section states four implications: what the existing standards already require and what this architecture operationalises beyond it, digital reporting, inspection measured precisely rather than declared, and verification economics.

\subsection{Standards Mapping}
\label{sec:standardsmap}

Table~\ref{tab:standardsmap} states, standard by standard, what is already required, what this architecture makes mechanical, what stays a matter of professional judgement no architecture should remove, and what becomes machine-inspectable as a consequence.

\begin{table}[htbp]
\centering
\footnotesize
\begin{tabular}{>{\raggedright\arraybackslash}p{1.9cm}>{\raggedright\arraybackslash}p{2.6cm}>{\raggedright\arraybackslash}p{2.6cm}>{\raggedright\arraybackslash}p{2.3cm}>{\raggedright\arraybackslash}p{2.3cm}}
\toprule
\textbf{Requirement} & \textbf{Standard requires} & \textbf{Discipline operationalises} & \textbf{Remains judgement} & \textbf{Becomes inspectable} \\
\midrule
Evidence sufficiency \& appropriateness & Sufficient, appropriate evidence, sources not ranked \citep{isa500} & Evidence grade (Table~\ref{tab:boundaryspec}) ranked mechanically per assertion class & Which grade suffices at which materiality & Whether the declared grade was actually met \\
\addlinespace
External confirmation & Auditor-controlled confirmation ordinarily more reliable \citep{isa505} & ``Route controlled by auditor'' as a checked precondition, not a heuristic & Whether a novel arrangement counts as auditor-controlled & Whether the route-control precondition held \\
\addlinespace
Management representations & Alone, not sufficient appropriate evidence \citep{isa580} & Representation-only evidence cannot satisfy a Verified postcondition alone & How much corroboration a specific representation needs & Whether a classification relied on representation alone \\
\addlinespace
Modified opinion & Qualified opinion or disclaimer for material scope limitations \citep{isa705} & \textsf{Undefined} plus an unresolved residual is the trigger, computed not judged & Whether a given \textsf{Undefined} region is material & Whether an unqualified opinion issued despite it \\
\addlinespace
Qualified/ disclaimer distinction & Pervasive-vs-not-pervasive branching \citep{cas705, pcaob3105} & Materiality threshold field operationalises the branching & The pervasiveness judgement itself & Whether the declared threshold was applied consistently \\
\addlinespace
Digital disclosure & Extensions must anchor to a standard taxonomy concept \citep{esma2019esef} & Boundary classification and residual status proposed as machine-readable tags (Section~\ref{sec:xbrl}) & Whether a specific anchoring is semantically apt & Whether every fact tagged audited carries a resolved classification \\
\bottomrule
\end{tabular}
\caption{Standards mapping: what each standard already requires, what Boundary Discipline operationalises beyond it, what remains professional judgement, and what becomes machine-inspectable. This table states an analysis, not a claim of regulatory adoption or legal sufficiency; no standard-setter has adopted any part of it.}
\label{tab:standardsmap}
\end{table}

\subsection{Digital Reporting: XBRL and ESEF}
\label{sec:xbrl}

Digital reporting formats are an application of this architecture, not formal-kernel functionality, and the distinction matters because XBRL's own guarantees are frequently overstated. Six things need to be kept apart. \emph{Syntactic validity} is whether a filing parses as well-formed XBRL. \emph{Taxonomy conformance} is whether every tag used exists in the governing taxonomy. \emph{Extension anchoring} is whether a preparer-authored extension element is linked, per the ESEF Regulatory Technical Standard, to a standard taxonomy concept of wider or narrower meaning \citep{esma2019esef}. \emph{Semantic correctness} is whether a tagged fact means what its tag says it means -- something no amount of syntactic or taxonomy checking establishes. \emph{Audit verification status} is this paper's own subject: whether the tagged fact is classified \textsf{Verified} under a declared boundary and validly admitted, or is \textsf{Undefined}, \textsf{Refuted}, or otherwise unresolved. \emph{Residual disclosure} is whether an unresolved fact's residual status is itself visible to a filing's consumer, rather than silently absent.

An XBRL filing that passes syntactic validity, taxonomy conformance, and anchoring establishes none of semantic correctness, audit verification status, or residual disclosure. Under Boundary Discipline, an unanchored or insufficiently justified extension is a boundary violation in exactly the sense defined here: automated validation confirms the filing is well-formed against a taxonomy the preparer itself authored, which is a syntactic check, and must surface as \textsf{Undefined} or an escalation rather than pass silently because the file parses. But anchoring, even when present and justified, is not itself a claim this paper makes about audit verification status; a fact can be perfectly anchored and entirely unverified.

A future machine-readable assurance extension, proposed here rather than claimed as an existing XBRL feature, could carry alongside a tagged fact: the identity and version of the boundary specification the fact was classified under, its classification state, its admission status, its residual identifier if any, and its disclosure or escalation status. This would let a filing's consumer distinguish, without access to the underlying working papers, a tagged fact backed by an admitted \textsf{Verified} classification from one carrying an open, disclosed residual -- which is the digital-reporting analogue of Remark~\ref{rem:success}: the consumer is told what is not known, not left to assume uniform assurance because every fact in the filing looks the same.

\subsection{Inspection Measures}
\label{sec:inspection}

Because boundary specifications and classification logs are machine-readable, inspection can evaluate precisely defined indicators instead of declared policy. Four measures are defined here; none of them establishes that the underlying classifications are substantively correct, only that they can be counted and compared.

\begin{definition}[Boundary coverage ratio]
\label{def:coverage}
For a set of material assertions $M$ and a declared boundary specification $B$, the boundary coverage ratio is
\[
\mathrm{coverage}(M, B) = \frac{|\{a \in M : \exists\, \text{a procedure in } B \text{ that can produce a determinate result for } a\}|}{|M|}.
\]
The denominator is every assertion the engagement's own materiality assessment has scoped as material (Section~\ref{sec:coq}), not every assertion in the specification; an assertion class entirely missing from $B$ counts against the numerator, as an assertion with no covering procedure, not as an omission from the denominator.
\end{definition}

\begin{definition}[Undefined-region materiality]
\label{def:undefmat}
For the subset $U \subseteq M$ classified \textsf{Undefined}, undefined-region materiality is reported as the aggregate amount $\sum_{a \in U} \mathrm{amount}(a)$ in the relevant reporting currency, and separately as that amount relative to the applicable financial-statement base (total assets, revenue, or the materiality benchmark the engagement has set). Amounts that are not additive across assertions, and purely qualitative materiality (a control deficiency, a going-concern indicator), must be disclosed as a separate qualitative list rather than forced into the sum; a single aggregate figure that silently drops non-additive items has reintroduced exactly the silent conversion this paper argues against, at the measurement layer.
\end{definition}

\begin{definition}[Residual assignment rate]
\label{def:assignrate}
For the set of emitted material residuals $R$, let $R_{\mathrm{assigned}} \subseteq R$ be those with an identified owner whose acceptance and capacity have both been recorded (Table~\ref{tab:rarschema}). The residual assignment rate is $|R_{\mathrm{assigned}}| / |R|$. A proposed-owner name with no recorded acceptance or capacity assessment does not count toward $R_{\mathrm{assigned}}$; Table~\ref{tab:rarentries}'s first illustrative entry is exactly this case and would not count as assigned under this definition.
\end{definition}

\begin{definition}[Orphan rate]
\label{def:orphanrate}
The orphan rate is $|R \setminus R_{\mathrm{assigned}}| / |R|$, the complement of the residual assignment rate over the same material residual set. Table~\ref{tab:rarentries}'s third illustrative entry is orphaned in this sense and contributes to this rate until an owner is proposed and accepted.
\end{definition}

These four measures support inspection of whether a firm's classification architecture actually eliminated implicit defaults; they do not support inspection of whether the underlying boundary specification, materiality thresholds, or classifications are themselves correct, which remains a matter of professional and regulatory judgement outside what any of these ratios can show. Illustrative requirements in the style of a boundary verification standard, including a requirement to report these four measures, are collected in Appendix~\ref{app:requirements}.

\subsection*{Verification economics}

Once residuals are explicit, they can be priced, insured, escalated, or disclosed. The value created by a verification procedure is the measurable difference between pre-verification and post-verification risk exposure, and explicit residual registers make that difference estimable in ways that opaque assurance does not. A fuller development of verification pricing is left to future work.

\section{Limitations and Future Work}
\label{sec:limitations}

The formal core covers classification logic, not audit substance. Boundary Discipline does not eliminate professional judgement, and it does not detect fraud; auditors still decide where boundaries are drawn, what evidence grades to assign, and what materiality thresholds to apply. The framework makes those decisions explicit, versioned, and inspectable, and it forbids their silent circumvention, which is a narrower and more defensible claim. The boundary specification of Table~\ref{tab:boundaryspec} is a proof-of-concept artefact requiring domain-specific validation before production use. The selected classification, admission, and residual-preservation requirements Table~\ref{tab:validation} reports are already machine-checked, generally, in the companion Coq/Rocq development described in Section~\ref{sec:coq}; what remains open is not their mechanisation but their application to a validated, production-grade boundary specification, and the substantive judgement, sampling sufficiency, valuation, and fraud detection, that the framework deliberately leaves outside its scope.

The evaluation this paper actually performs is design-science evaluation, not empirical validation, and the two must not be conflated. This paper does not claim practitioner validation, field deployment, or measured effects on deficiency rates, restatement frequency, exception closure, residual visibility, or modified-opinion decisions -- none of these questions have data behind them here. A future evaluation programme, appropriate to a design-science artefact at this stage, would need: structured expert interviews with auditors and AIS specialists on the boundary specification's construct validity; controlled workflow studies comparing disciplined and undisciplined exception handling on the same transaction population; field implementation at a firm willing to run the architecture alongside its existing methodology; inspection reproducibility, whether two independent inspectors compute the same Section~\ref{sec:inspection} measures from the same classification logs; and, only once the above exist, archival study of whether default-free classification measurably changes deficiency rates, restatement frequency, exception closure rates, residual visibility, or the incidence of modified opinions. None of this paper's claims depend on that programme having been run.

More specifically, and stated here rather than only in the companion repository: this development does not establish that supplied evidence is genuine -- a claim packet's preconditions are asserted, not independently verified, by whatever process constructs it. It does not establish that supplied materiality judgements are correct; materiality is a caller-supplied threshold this architecture applies, not one it derives. It does not establish that distinct registered process identities are free of collusion or organisational dependence -- Section~\ref{sec:pipeline}'s qualification on the admission verifier applies throughout. It does not establish that the identifiers a caller supplies for processes or proposals are unique, and a caller that reuses one can produce a certificate that validates against a proposal other than the one it was drawn up for. It does not establish that the OCaml extraction and compilation step is itself machine-checked, rather than trusted in the way any compiled artefact is. And it does not establish that a residual, once recorded, is ever formally closed or resolved by the kernel; the register only ever grows. See the companion repository's \texttt{NON\_CLAIMS.md} for the itemised list this paragraph summarises.

Some audit settings also show that documentary evidence can be proactively manufactured, so the boundary problem is not only missing evidence but misplaced authority in the evidence hierarchy; a fuller treatment belongs to a critical accounting venue. Finally, AI-assisted audit systems make the problem addressed here more urgent rather than different in kind. Generated conclusions should be treated as proposed classifications requiring independent admission, never as verified outputs, and the admission-separation property (Table~\ref{tab:validation}) was designed with exactly this extension in view. Its development is reserved for separate work.

\section{Conclusion}
\label{sec:conclusion}

This paper asked how an accounting information system should be designed so that verification failure remains a durable, consequence-bearing state from evidence intake through audit reporting, residual-risk assignment, inspection, and digital disclosure. The answer is architectural, not merely definitional: verification failure must remain represented through every layer of the system, from the boundary specification that computes \textsf{Undefined} through the Verification Balance Sheet that states what a procedure actually established, to the Risk Assignment Register that keeps an unresolved residual visible after the workflow that flagged it has closed. Unresolved material assertions must constrain audit consequence, not merely be recorded somewhere for later; residuals require an owner with recorded acceptance and capacity, or an orphan flag that reaches an escalation authority, not a name in a field no one reviews. Boundary Discipline makes these conditions inspectable -- through the boundary coverage ratio, undefined-region materiality, residual assignment rate, and orphan rate of Section~\ref{sec:inspection} -- rather than merely asserted as policy.

The architecture does not eliminate professional judgement, audit risk, or fraud; its claim is narrower and harder to dismiss. What it establishes beyond that claim, as supporting evaluation evidence rather than as this paper's own technical contribution, is that a companion, independently machine-checked kernel demonstrates the selected classification, admission, and preservation requirements this architecture depends on can be implemented, together, without an affirmative default. The technical resources for this discipline already exist in audit theory, in accounting information systems research, in database semantics, and in runtime verification. The contribution of this paper is the design-science architecture that joins them into something an engagement can build, and an inspector can measure.

% --------------------------------------------------------------------
% Reference list formatted per IJDAR's Reference Style page (APA 7th
% edition; https://www.ijdar.org/reference-style), arranged alphabetically.
% DOIs below were resolved against publisher and Crossref-indexed records.
% Ten entries (cas705, esma2019esef, esma2020, isa500, isa505, isa580,
% isa705, pcaob3105, vasarhelyi1991, wambach2021) carry no DOI because
% none has been assigned by the publisher (regulatory standards,
% official/grey literature, and an unreleased Bundestag investigation
% report whose text is only available via press republication); verify
% this remains current before submission rather than inserting a
% placeholder or guessed DOI. cas705 and pcaob3105 previously lacked
% even the "No DOI assigned" annotation the other entries carry --
% fixed for consistency, not newly discovered as missing a DOI. The
% isa500/505/580/705 titles and "(Revised and Redrafted)" designations
% follow the 2016 IAASB Handbook; wambach2021's submission and
% publication dates were corroborated via press reporting
% (Handelsblatt), not the Bundestag itself, which has not released the
% report.
% --------------------------------------------------------------------
\begin{thebibliography}{99}

\bibitem[Alles et al.(2006)]{alles2006}
Alles, M.~G., Brennan, G., Kogan, A., \& Vasarhelyi, M.~A. (2006).
Continuous monitoring of business process controls: A pilot implementation of a continuous auditing system at Siemens.
\emph{International Journal of Accounting Information Systems}, 7(2), 137--161. \url{https://doi.org/10.1016/j.accinf.2005.10.004}

\bibitem[Austin et al.(2021)]{austin2021}
Austin, A.~A., Carpenter, T.~D., Christ, M.~H., \& Nielson, C.~S. (2021).
The data analytics journey: Interactions among auditors, managers, regulation, and technology.
\emph{Contemporary Accounting Research}, 38(3), 1888--1924. \url{https://doi.org/10.1111/1911-3846.12680}

\bibitem[Bauer et al.(2010)]{bauer2010}
Bauer, A., Leucker, M., \& Schallhart, C. (2010).
Comparing LTL semantics for runtime verification.
\emph{Journal of Logic and Computation}, 20(3), 651--674. \url{https://doi.org/10.1093/logcom/exn075}

\bibitem[Bauer et al.(2011)]{bauer2011}
Bauer, A., Leucker, M., \& Schallhart, C. (2011).
Runtime verification for LTL and TLTL.
\emph{ACM Transactions on Software Engineering and Methodology}, 20(4), Article 14. \url{https://doi.org/10.1145/2000799.2000800}

\bibitem[Beerbaum(2020)]{beerbaum2020}
Beerbaum, D. (2020).
The future of audit after the Wirecard accounting scandal: Consideration of the `Wambach report' (Working paper).
\emph{Aalto University School of Business, Department of Accounting}. \url{https://doi.org/10.2139/ssrn.4164941}
Secondary commentary on \citet{wambach2021}, not the investigation itself; cited here for its discussion, with the investigation cited separately for its findings.

\bibitem[CPA Canada(2016)]{cas705}
Chartered Professional Accountants of Canada. (2016).
\emph{CAS 705: Modifications to the opinion in the independent auditor's report}. CPA Canada Handbook -- Assurance. No DOI assigned.

\bibitem[Codd(1979)]{codd1979}
Codd, E.~F. (1979).
Extending the database relational model to capture more meaning.
\emph{ACM Transactions on Database Systems}, 4(4), 397--434. \url{https://doi.org/10.1145/320107.320109}

\bibitem[European Securities and Markets Authority(2019)]{esma2019esef}
European Securities and Markets Authority. (2019).
\emph{Commission Delegated Regulation (EU) 2019/815 of 17 December 2018 supplementing Directive 2004/109/EC with regard to regulatory technical standards on the specification of a single electronic reporting format} (ESEF Regulation). \emph{Official Journal of the European Union}, L 143. No DOI assigned.

\bibitem[ESMA(2020)]{esma2020}
European Securities and Markets Authority. (2020).
\emph{Fast track peer review on the application of the guidelines on the enforcement of financial information by BaFin and FREP in the context of Wirecard} (Report No. ESMA42-111-5349). No DOI assigned.

\bibitem[Habib(2013)]{habib2013}
Habib, A. (2013).
A meta-analysis of the determinants of modified audit opinion decisions.
\emph{Managerial Auditing Journal}, 28(3), 184--216. \url{https://doi.org/10.1108/02686901311304349}

\bibitem[Hevner et al.(2004)]{hevner2004}
Hevner, A.~R., March, S.~T., Park, J., \& Ram, S. (2004).
Design science in information systems research.
\emph{MIS Quarterly}, 28(1), 75--105. \url{https://doi.org/10.2307/25148625}

\bibitem[Hurtt(2010)]{hurtt2010}
Hurtt, R.~K. (2010).
Development of a scale to measure professional skepticism.
\emph{Auditing: A Journal of Practice \& Theory}, 29(1), 149--171. \url{https://doi.org/10.2308/aud.2010.29.1.149}

\bibitem[IAASB(2016a)]{isa500}
International Auditing and Assurance Standards Board. (2016).
\emph{International Standard on Auditing 500: Audit evidence}. IFAC. No DOI assigned.

\bibitem[IAASB(2016b)]{isa505}
International Auditing and Assurance Standards Board. (2016).
\emph{International Standard on Auditing 505 (Revised and Redrafted): External confirmations}. IFAC. No DOI assigned.

\bibitem[IAASB(2016c)]{isa580}
International Auditing and Assurance Standards Board. (2016).
\emph{International Standard on Auditing 580 (Revised and Redrafted): Written representations}. IFAC. No DOI assigned.

\bibitem[IAASB(2016d)]{isa705}
International Auditing and Assurance Standards Board. (2016).
\emph{International Standard on Auditing 705 (Revised): Modifications to the opinion in the independent auditor's report}. IFAC. No DOI assigned.

\bibitem[Issa and Kogan(2014)]{issa2016}
Issa, H., \& Kogan, A. (2014).
A predictive ordered logistic regression model as a tool for quality review of control risk assessments.
\emph{Journal of Information Systems}, 28(2), 209--229. \url{https://doi.org/10.2308/isys-50808}

\bibitem[Knechel et al.(2013)]{knechel2013}
Knechel, W.~R., Krishnan, G.~V., Pevzner, M., Shefchik, L.~B., \& Velury, U.~K. (2013).
Audit quality: Insights from the academic literature.
\emph{Auditing: A Journal of Practice \& Theory}, 32(Supplement 1), 385--421. \url{https://doi.org/10.2308/ajpt-50350}

\bibitem[Kogan et al.(1999)]{kogan1999}
Kogan, A., Sudit, E.~F., \& Vasarhelyi, M.~A. (1999).
Continuous online auditing: A program of research.
\emph{Journal of Information Systems}, 13(2), 87--103. \url{https://doi.org/10.2308/jis.1999.13.2.87}

\bibitem[Leucker and Schallhart(2009)]{leucker2009}
Leucker, M., \& Schallhart, C. (2009).
A brief account of runtime verification.
\emph{The Journal of Logic and Algebraic Programming}, 78(5), 293--303. \url{https://doi.org/10.1016/j.jlap.2008.08.004}

\bibitem[Mautz and Sharaf(1961)]{mautz1961}
Mautz, R.~K., \& Sharaf, H.~A. (1961).
\emph{The philosophy of auditing}. American Accounting Association.

\bibitem[McCrum(2022)]{mccrum2022}
McCrum, D. (2022).
\emph{Money men: A hot startup, a billion dollar fraud, a fight for the truth}. Bantam Press.

\bibitem[Nelson(2009)]{nelson2009}
Nelson, M.~W. (2009).
A model and literature review of professional skepticism in auditing.
\emph{Auditing: A Journal of Practice \& Theory}, 28(2), 1--34. \url{https://doi.org/10.2308/aud.2009.28.2.1}

\bibitem[PCAOB(2017)]{pcaob3105}
Public Company Accounting Oversight Board. (2017).
\emph{AS 3105: Departures from unqualified opinions and other reporting circumstances} (PCAOB Release No. 2017-001). No DOI assigned.

\bibitem[Peffers et al.(2007)]{peffers2007}
Peffers, K., Tuunanen, T., Rothenberger, M.~A., \& Chatterjee, S. (2007).
A design science research methodology for information systems research.
\emph{Journal of Management Information Systems}, 24(3), 45--77. \url{https://doi.org/10.2753/MIS0742-1222240302}

\bibitem[Svanberg et al.(2025)]{svanberg2025}
Svanberg, J., Khoshkhou, N.~F., \"{O}hman, P., \& L\"{o}vqvist, D. (2025).
Addressing the exception prioritization problem in continuous auditing systems with thresholding.
\emph{Intelligent Systems in Accounting, Finance and Management}, 32(1), e70022. \url{https://doi.org/10.1002/isaf.70022}

\bibitem[Vasarhelyi and Halper(1991)]{vasarhelyi1991}
Vasarhelyi, M.~A., \& Halper, F.~B. (1991).
The continuous audit of online systems.
\emph{Auditing: A Journal of Practice \& Theory}, 10(1), 110--125. No DOI assigned.

\bibitem[Wambach(2021)]{wambach2021}
Wambach, M. (2021).
\emph{Sonderpr\"{u}fungsbericht: Bericht \"{u}ber die Ergebnisse der Untersuchung gem\"{a}\ss{} Beweisbeschluss SV-4 des 3. Untersuchungsausschusses der 19. Wahlperiode des Deutschen Bundestages} [Special investigation report to the German Bundestag's Wirecard inquiry committee]. R\"{o}dl \& Partner. Submitted 16 April 2021; not released by the Bundestag, text published by Handelsblatt, 12 November 2021. No DOI assigned.

\end{thebibliography}

\appendix

\section{Glossary of Terms}
\label{app:glossary}

\noindent\textbf{Admission.} The act by which a distinct, registered verifier identity accepts a proposed classification as binding for audit consequence.

\medskip
\noindent\textbf{Boundary.} The demarcation between regions where a verification procedure produces valid results and regions where it does not.

\medskip
\noindent\textbf{Fiat Ground.} The foundational assumptions a verification system takes as given, which cannot themselves be verified within the system.

\medskip
\noindent\textbf{Implicit default.} A classification rule that applies when explicit procedures fail, typically producing affirmative outputs without documentation.

\medskip
\noindent\textbf{Residual.} What remains unverified after a verification procedure completes, including both known limitations and structural blind spots.

\medskip
\noindent\textbf{Risk assignment.} The process of identifying who bears responsibility for losses if a verification residual materialises as actual error.

\medskip
\noindent\textbf{Silent conversion.} The operational treatment of a verification failure as acceptance without an explicit \textsf{Undefined}, modified, blocked, or escalated classification.

\medskip
\noindent\textbf{Trust Chain.} The sequence of delegations connecting a verification procedure's output to its fiat ground.

\medskip
\noindent\textbf{Undefined.} The classification result when an input falls outside all defined verification boundaries.

\medskip
\noindent\textbf{Verification failure.} The condition in which a system cannot determine whether a claim is verified or refuted within the evidence, authority, and observation boundaries available to it.

\section{Illustrative Boundary Verification Requirements}
\label{app:requirements}

The following requirements illustrate how Boundary Discipline could be codified in an auditing standard, an inspection protocol, or an internal firm methodology. They are illustrative rather than proposed regulatory text.

\begin{enumerate}
\item Auditors must define explicit verification boundaries for all material assertion classes, specifying applicable methods, jurisdictions, evidence grades, and materiality thresholds, documented in a machine-readable format subject to version control and audit trail requirements.
\item Audit methodology must produce \textsf{Undefined} rather than affirmative classifications for assertions outside defined boundaries. Management representations cannot substitute for primary evidence above materiality thresholds, scope limitations affecting material items preclude unqualified opinions, and documentation from unverified intermediaries cannot establish asset existence.
\item For engagements above a size threshold reflecting systemic importance and fee capacity, the classification logic mapping evidence to conclusions must be specified in a formal language amenable to machine verification, proven free of implicit defaults using a proof assistant, and extracted to executable code used in actual audit procedures.
\item Audit reports must include a Risk Assignment Register (Table~\ref{tab:rarschema}) identifying each verification residual, its proposed owner, whether that owner's acceptance and capacity have been recorded, and any orphaned residuals -- an owner name alone, without recorded acceptance and capacity, must not count as completed assignment.
\item Inspections and internal quality-control reviews must evaluate the boundary coverage ratio, undefined-region materiality, residual assignment rate, and orphan rate (Section~\ref{sec:inspection}) derived from the boundary specifications and classification logs.
\end{enumerate}

Requirements of type (1) through (4) are compliance requirements that firms must satisfy; requirement (5) comprises inspectable indicators that make compliance measurable. The distinction prevents paper compliance, meaning a boundary regime declared in policy but not enforced in the operative classification procedure.

\end{document}
